% ============================================================
%  Statische Geometrie und SI-Projektion
%  Warum natuerliche Einheiten keine fraktalen Korrekturen brauchen
%  Kapitelinhalt DE
%  Quellenbelegt: Dok. 011, 012, 013, 014, 015, 041, 044, 116, 133, 193, 258
% ============================================================

\title{Dok.~261 — Statische Geometrie und SI-Projektion:\\
Warum natürliche Einheiten keine fraktalen Korrekturen brauchen}
\author{Johann Pascher}
\date{Mai 2026}
\maketitle

\begin{center}
\textit{Kontext: Dok.~011 (Feinstruktur), Dok.~013/014 (SI \& natürliche Einheiten),
Dok.~116 (Koide), Dok.~133 (fraktale Korrektur)} \\[2pt]
\textit{FFGFT-Archiv: DOI 10.5281/zenodo.20117635}
\end{center}

\tableofcontents
\newpage

%=============================================================
\section{Einordnung}
%=============================================================

Die SI-Reform von 2019 hat vier Konstanten ($c$, $\hbar$, $e$, $k_B$) auf exakte
Werte fixiert und damit die meisten dimensionsbehafteten Größen zu reinen
Einheitendefinitionen gemacht. Was als \emph{echter} empirischer Input übrig
bleibt, sind die dimensionslosen Konstanten -- allen voran die
Feinstrukturkonstante $\alpha$, deren Wert seit der Reform weiterhin gemessen
wird ($\mu_0$, $\varepsilon_0$ und $\alpha$ blieben Messgrößen). FFGFT leitet
auch diese dimensionslosen Inputs aus dem einzigen Parameter
$\xipar = \tfrac{4}{3}\times 10^{-4}$ ab.

Dieses Dokument stellt die Umrechnungssystematik kompakt zusammen
(Tabellen~\ref{tab:si-fixiert}--\ref{tab:xi-inputs}) und begründet anschließend
einen zentralen Punkt: \emph{solange man in natürlichen Einheiten und reinen
Verhältnissen bleibt, ist die Beschreibung statisch und relativ und erfordert
keine fraktalen Korrekturen.} Die fraktalen Korrekturen treten erst beim Übergang
in absolute SI-Werte auf.

%=============================================================
\section{Die durch die SI-Reform fixierten Konstanten}
%=============================================================

\begin{table}[htbp]
  \centering
  \caption{Die durch die SI-Reform 2019 fixierten Konstanten (die ``Maschinerie'').
           Sie tragen keine Physik, nur Maßstab, und werden in natürlichen
           Einheiten zu~$1$.}
  \label{tab:si-fixiert}
  \begin{tabular}{llll}
    \toprule
    Konstante & Wert (exakt seit 2019) & definiert & nat.\ Einheiten \\
    \midrule
    $c$     & $299\,792\,458$ m/s                          & Meter   & $=1$ \\
    $\hbar$ & $1{,}054\,571\,817\ldots\times10^{-34}$ J\,s & kg      & $=1$ \\
    $e$     & $1{,}602\,176\,634\times10^{-19}$ C           & Ampere  & Ladungseinheit \\
    $k_B$   & $1{,}380\,649\times10^{-23}$ J/K              & Kelvin  & $=1$ \\
    \bottomrule
  \end{tabular}
\end{table}

%=============================================================
\section{Größenumrechnung in natürlichen Einheiten}
%=============================================================

\begin{table}[htbp]
  \centering
  \caption{Größenumrechnung in natürlichen Einheiten ($\hbar=c=k_B=1$):
           alle dimensionsbehafteten Größen werden zu Potenzen der Energie.}
  \label{tab:nat-umrechnung}
  {\small
  \begin{tabular}{lccl}
    \toprule
    Größe & SI-Dimension & $[\hbar=c=k_B=1]$ & SI-Rückrechnung \\
    \midrule
    Energie       & J        & $E$           & Basis \\
    Masse         & kg       & $E$           & $\times\,S_{\mathrm{T0}}$ \ (über $c^{-2}$) \\
    Impuls        & kg\,m/s  & $E$           & $\times\,c^{-1}$ \\
    Länge         & m        & $E^{-1}$      & $\ell[\mathrm{fm}]=197{,}327/E[\mathrm{MeV}]$ \\
    Zeit          & s        & $E^{-1}$      & $t[\mathrm{s}]=6{,}582\times10^{-22}/E[\mathrm{MeV}]$ \\
    Temperatur    & K        & $E$           & $T[\mathrm{K}]=E[\mathrm{eV}]/(8{,}617\times10^{-5})$ \\
    Ladung (Gauß) & C        & dimensionslos & $e=\sqrt{\alpha}\approx0{,}0854$ \\
    \bottomrule
  \end{tabular}}

  \medskip
  \footnotesize
  Die SI-Rückrechnung läuft über die \emph{eine} Massenskala
  $S_{\mathrm{T0}} = 1\,\mathrm{MeV}/c^2 = 1{,}782662\times10^{-30}$~kg (Dok.~013/014);
  der Längenfaktor ist $\hbar c/S_{\mathrm{T0}} = 1{,}973\times10^{-13}$~m, also die
  $197$\,MeV\,fm aus der $\hbar c=1$-Spalte in T0-Sprache.\par\smallskip
  In T0-Einheiten wird die Ladungseinheit so gewählt
  ($e=\sqrt{4\pi\varepsilon_0\hbar c}$), dass der \emph{geschriebene} Wert
  $\alpha=1$ ist (Dok.~044) -- eine Konvention, kein physikalischer Eingriff.
  Die physikalische Invariante $\alpha\approx1/137$ wird aus $\xipar$ rekonstruiert
  (Tab.~\ref{tab:xi-inputs}).
\end{table}

%=============================================================
\section{Die dimensionslosen Inputs und ihre Herleitung aus $\xipar$}
%=============================================================

\begin{table}[htbp]
  \centering
  \caption{Die überlebenden dimensionslosen Inputs der Physik und ihre
           Herleitung aus dem einzigen Parameter $\xipar=\tfrac{4}{3}\times10^{-4}$.
           Sämtliche Relationen quellenbelegt.}
  \label{tab:xi-inputs}
  {\footnotesize\setlength{\tabcolsep}{4pt}
  \begin{tabular}{lccl}
    \toprule
    Invariante & Standardphysik & FFGFT-Relation (aus $\xipar$) & Quelle \\
    \midrule
    $\xipar$           & ---            & $=\tfrac{4}{3}\times10^{-4}=4/30000$ & Dok.~013 \\
    $\alpha$ (SI)      & $1/137{,}036$  & $\xipar\,(\Ezero/1\,\mathrm{MeV})^2$\,$^{\dagger}$ & Dok.~011 \\
    Leptonmassen       & gemessen       & $m=r\,\xipar^{\,p}\,v$ & Dok.~116, 006 \\
    \quad Exponenten   &                & $(p_e,p_\mu,p_\tau)=(\tfrac{3}{2},\,1,\,\tfrac{2}{3})$ & Dok.~116 \\
    \quad Verhältnisse &                & $(r_e,r_\mu,r_\tau)=(\tfrac{4}{3},\,\tfrac{16}{5},\,\tfrac{25}{9})$ & Dok.~116 \\
    Koide $Q$          & $0{,}66666$    & $=\tfrac{2}{3}$ \ (Folge der $\xipar$-Exponenten) & Dok.~116, 258 \\
    $\sin^2\theta_W$   & $0{,}2312$     & $\dfrac{1-\sqrt{1-4\alpha_W}}{4}$,\ \ $\alpha_W=\sqrt{\xipar}$ & Dok.~041 \\[0.6em]
    $\theta_{13}$      & $8{,}57^\circ$ & $\arcsin\!\big(\xipar^{1/3}\big)$ & Dok.~041 \\
    $G$                & gemessen       & $\dfrac{\xipar^2}{4m_e}\,\Kfrak$,\ \ $\Kfrak=1-100\xipar\approx0{,}9867$ & Dok.~012, 133 \\
    \bottomrule
  \end{tabular}}

  \medskip
  \footnotesize
  $^{\dagger}$\,$\alpha=\xipar\,(\Ezero/1\,\mathrm{MeV})^2$ ist die \emph{Leitordnung},
  mit $\Ezero=\sqrt{m_e\,m_\mu}=7{,}398$\,MeV.
  Die CODATA-Präzision wird über das fraktale Korrekturprodukt
  $\displaystyle\prod_{n=1}^{137}\!\big(1+\delta_n\,\xipar\,(4/3)^{\,n-1}\big)$
  erreicht. Diese Korrekturen sind \textbf{keine freien Annahmen}:
  $\Kfrak=1-100\xipar$ ist über den Skalenfluss von der Planck- zur
  elektroschwachen Skala eindeutig fixiert ($D_f^{\,\mathrm{Raum}}=3-\xipar$,
  Dok.~133, ``$D_f$ ist eindeutig bestimmt -- nicht frei wählbar''), und die
  rationalen Leiter-Verhältnisse $r_i$ sind $\alpha$-geometrische Invarianten
  (Dok.~258), strukturell unzugänglich für nachträgliche Anpassung
  (Nichtabschluss-Theorem, Dok.~193).
\end{table}

%=============================================================
\section{Statische Geometrie: warum natürliche Einheiten korrekturfrei sind}
%=============================================================

Solange die Physik in natürlichen Einheiten und reinen Verhältnissen ausgedrückt
wird, ist ihre Beschreibung \emph{statisch} und \emph{relational}: Es treten nur
dimensionslose Verhältnisse auf, keine absoluten Anker und keine
Skalenabhängigkeit. Auf dieser Ebene sind die FFGFT-Größen exakte geometrische
Invarianten von $\xipar$ -- und es ist keine fraktale Korrektur erforderlich.

\begin{prinzip}
\textbf{Prinzip.} In natürlichen Einheiten ist FFGFT rein relational. Die
dimensionslosen Strukturrelationen -- Massenverhältnisse, Mischungswinkel, die
Koide-Größe -- sind statische geometrische Invarianten von $\xipar$. Sie
enthalten keine fraktalen Korrekturen.
\end{prinzip}

Das deutlichste Beispiel ist die Koide-Größe: $Q=\tfrac{2}{3}$ folgt exakt aus den
Exponenten $(p_e,p_\mu,p_\tau)=(\tfrac{3}{2},1,\tfrac{2}{3})$ und den
Verhältnissen $(r_e,r_\mu,r_\tau)=(\tfrac{4}{3},\tfrac{16}{5},\tfrac{25}{9})$ --
ohne zusätzliche Parameter und ohne fraktale Korrektur (Dok.~116). Ebenso sind die
Leiter-Verhältnisse $r_i$, die Exponenten $p_i$, die Weinberg-Struktur und
$\theta_{13}$ exakte Funktionen von $\xipar$ allein. Keine dieser Relationen
benötigt einen absoluten Maßstab, ein $\Ezero$ in MeV oder einen
Umrechnungsfaktor.

Wo entstehen dann die fraktalen Korrekturen? Erst an der \emph{SI-Projektion} --
dort, wo eine absolute Skala festgelegt wird. $\Kfrak=1-100\xipar$ ist nicht eine
Korrektur an der Geometrie, sondern die kumulative Wirkung des Skalenflusses von
der Planck- zur elektroschwachen Skala (Dok.~133): eine Eigenschaft der Abbildung
der Geometrie auf ein absolutes Energie-Lineal. Genauso erscheint das fraktale
Korrekturprodukt $\prod_n(1+\delta_n\,\xipar\,(4/3)^{n-1})$ erst, wenn die führende
Beziehung $\alpha=\xipar(\Ezero/1\,\mathrm{MeV})^2$ auf CODATA-Präzision in
absoluten SI-Werten gebracht wird.

Wichtig zur Abgrenzung: Diese Projektionsfaktoren sind dennoch \emph{keine} freien
Parameter. $\Kfrak$ ist über den Skalenfluss eindeutig fixiert (Dok.~133), die
rationalen Verhältnisse sind $\alpha$-geometrische Invarianten (Dok.~258), und das
Nichtabschluss-Theorem (Dok.~193) sperrt jede nachträgliche Anpassung. Sie sind
abgeleitete Projektionsfaktoren, nicht angenommene Fits.

\begin{schlussfolgerung}
\textbf{Schlussfolgerung.} Die fraktalen Korrekturen sind der Preis der
Absolut-Projektion, nicht ein Defekt der Geometrie. Solange man in natürlichen
Einheiten und Verhältnissen bleibt, ist die Beschreibung statisch, relativ und
korrekturfrei; die Korrekturen treten ausschließlich beim Übergang in absolute
SI-Werte auf -- und sind selbst dort eindeutig aus $\xipar$ und der
Skalenstruktur bestimmt.
\end{schlussfolgerung}
